\documentclass[11pt,a4paper]{article}
\usepackage[utf8]{inputenc}
\usepackage[french]{babel}
\usepackage[T1]{fontenc}
\usepackage{amsmath, amssymb, amsthm}
\usepackage{geometry}
\geometry{margin=2.5cm}

\title{Modélisation du Problème de Covoiturage par un problème de Matching avec Coût}
\author{NAOUSSI DJOUMSSI Viery}

\begin{document}

\maketitle

\section{Introduction}
Ce document définit formellement le problème de covoiturage par un problème de matching "many-to-one" car chaque passagé est affecté à une voiture unique et une voiture peut acceuillir plusieurs passager.
\section{Définition des Concepts et Variables}

\subsection{Les Agents et le Matching Many-to-One}
Le problème est modélisé comme un \textbf{matching Many-to-One} (plusieurs-à-un) de l'ensemble $P$ vers l'ensemble $D'$ où :
\begin{itemize}
    \item $P = \{j_1, ..., j_n\}$ l'ensemble des passagers.
    \item $D' = D \cup F$ l'ensemble étendu des véhicules, où $D$ sont les conducteurs réels et $F = \{f_1, ..., f_n\}$ sont les \textbf{voitures fantômes} (une par passager).
\end{itemize}

\subsection{Variables de Décision}
On définit la variable binaire $x_{i,j}$ pour tout $i \in D'$ et $j \in P$ :
\begin{equation}
    x_{i,j} = \begin{cases}
    1 & \text{si le passager } j \text{ est affecté au véhicule } i \\
    0 & \text{sinon}
    \end{cases}
\end{equation}

\subsection{La Matrice de Coût $C_{i,j}$}
$C_{i,j}$ représente le coût généralisé induit par l'affectation du passager $j$ au véhicule $i$. Sa valeur dépend de la nature du véhicule :
\begin{equation}
    C_{i,j} = \begin{cases}
    C^C_{i,j} & \text{si } i \in D \text{ (Covoiturage)} \\
    C^M_j & \text{si } i \in F \text{ et } i \text{ correspond à } j \text{ (Marche)} \\
    +\infty & \text{si } i \in F \text{ et } i \text{ ne correspond pas à } j
    \end{cases}
\end{equation}

\subsection{Désignations Temporelles}
Pour chaque passager $j$, nous définissons les variables temporelles suivantes :
\begin{itemize}
    \item $t^h_j$ : l'heure à laquelle le passager quitte la maison.
    \item $t^{c_i}_{m,j}$ : l'heure à laquelle le passager rencontre le conducteur $i$.
    \item $tt^c_{i,j}$ : le temps total de trajet en covoiturage avec le conducteur $i$.
    \item $tt^M_j$ : le temps total de marche si le passager décide de marcher jusqu'à sa destination.
    \item $t^*_j$ : l'heure d'arrivée désirée du passager.
    \item $t_{a,j}$ : l'heure d'arrivée réelle à destination.
\end{itemize}

\subsection{Définition Analytique des Coûts}
Le coût est une fonction du temps passé en transport et de la ponctualité à l'arrivée.
\begin{itemize}
    \item \textbf{Coût de covoiturage} ($C^C_{i,j}$) :
    \begin{equation}
        C^C_{i,j} = \alpha^c \cdot tt^c_{i,j} + SDC_j + \alpha^M \cdot tt^M_{access,j}
    \end{equation}
    Où $tt^M_{access,j}$ est le temps de marche résiduel (accès/egress).
    \item \textbf{Coût de marche} ($C^M_j$) :
    \begin{equation}
        C^M_j = \alpha^M \cdot tt^M_j + SDC_j
    \end{equation}
    Nous considérons que $tt^c_{i,j} < tt^{M}_j$ et $\alpha^c < \alpha^M$.
\end{itemize}

\section{Concepts de Compatibilité}

\subsection{Compatibilité Temporelle}
Le passager $j$ ne peut être pris en charge par le conducteur $i$ que si l'heure de rencontre $t^{c_i}_{m,j}$ respecte son heure de départ au plus tôt : $t^{c_i}_{m,j} \ge t^h_j$.

\subsection{Compatibilité Sociale}
Critère binaire ou probabiliste imposant une structure de clique : deux passagers affectés au même véhicule $i$ doivent être compatibles entre eux.

\section{Le Schedule Delay Cost (SDC) : Analyse des Formes Possibles}

Le $SDC$ mesure l'insatisfaction liée à la différence entre l'heure d'arrivée réelle ($t_a$) et l'heure souhaitée ($t^*$). Plusieurs modélisations sont possibles selon le comportement recherché :

\begin{enumerate}
    \item \textbf{Modèle Linéaire Symétrique} :
    \begin{equation}
        SDC(t_a) = \alpha \cdot |t^* - t_a|
    \end{equation}
    L'avance et le retard sont pénalisés de la même manière.

    \item \textbf{Modèle Quadratique} :
    \begin{equation}
        SDC(t_a) = \alpha \cdot (t^* - t_a)^2
    \end{equation}
    Pénalise très lourdement les grands écarts.
    \item \textbf{Modèle Asymétrique Linéaire (Small, 1982)} :
    \begin{equation}
        SDC(t_a) = \beta \cdot (t^* - t_a)^+ + \gamma \cdot (t_a - t^*)^+
    \end{equation}
    On considère généralement $\gamma > \alpha > \beta$ car être en retard est plus pénalisant (perte de productivité, rendez-vous manqué) qu'être en avance (temps d'attente passif).

    \item \textbf{Modèle avec Pénalité Fixe de Retard} :
    \begin{equation}
        SDC(t_a) = \beta \cdot (t^* - t_a)^+ + \gamma \cdot (t_a - t^*)^+ + \theta \cdot \mathbb{1}_{t_a > t^*}
    \end{equation}
    Où $\theta$ est un coût fixe subi dès que l'heure d'arrivée dépasse $t^*$, indépendamment de la durée du retard.
\end{enumerate}

\section{Modélisation Mathématique (ILP)}

\subsection{Fonction Objectif}
\begin{equation}
    \min Z = \sum_{i \in D'} \sum_{j \in P} x_{i,j} \cdot C_{i,j}
\end{equation}

\subsection{Contraintes}
\begin{enumerate}
    \item \textbf{Unicité} : $\forall j \in P, \sum_{i \in D'} x_{i,j} = 1$.
    \item \textbf{Capacité} : $\forall i \in D, \sum_{j \in P} x_{i,j} \le K$.
    \item \textbf{Clique} : $\forall i \in D, \forall (j_1, j_2) \in P^2$, si $j_1, j_2$ incompatibles, alors $x_{i,j_1} + x_{i,j_2} \le 1$.
    \item \textbf{Compatibilité Temporelle} : $t^{c_i}_{m, j} \geq t^{h}_j$.
    \item \textbf{Rationnalité individuelle} : si un passager $j$ est associé à un conducteur $i$ alors $C_j^M \geq C_{i, j}^C$.
\end{enumerate}

\end{document}
